Este informe detalla el estudio práctico de dos problemas clásicos en las ciencias de la computación, uno de ellos siendo el ordenamiento de arreglos unidimensionales y la multiplicación de matrices cuadradas. Para el primer problema, se evalúa el rendimiento de los algoritmos \textit{Standard Sort}, \textit{Merge Sort}, \textit{Quick Sort} y \textit{Patience Sort} bajo distintas distribuciones de datos iniciales. Para el segundo, se contrastan los enfoques del algoritmo estándar (\textit{Naive}) y el método de división y conquista de \textit{Strassen}. El principal objetivo del análisis experimental es relacionar el comportamiento práctico de cada algoritmo, midiendo su tiempo de ejecución y su consumo de memoria RAM, con su complejidad temporal y espacial teórica, validando su eficiencia en un entorno real.